% BHU PYQ -- Manually transcribed & error-corrected
% Paper: GLB-101: Elementary Physical and Structural Geology
% Exam: B.Sc. (Hons.) Semester I Examination, 2023-24

\documentclass[11pt,a4paper]{article}
\usepackage[a4paper,top=2.5cm,bottom=2.5cm,left=2.8cm,right=2.8cm,headheight=14pt]{geometry}
\usepackage{amsmath,amssymb}
\usepackage[T1]{fontenc}
\usepackage[utf8]{inputenc}
\usepackage{lmodern,microtype,fancyhdr,enumitem,hyperref,lastpage,parskip}

\hypersetup{colorlinks=true,urlcolor=black,linkcolor=black,
  pdftitle={GLB-101 Elementary Physical and Structural Geology -- BHU B.Sc. Sem I 2023-24}}
\pagestyle{fancy}\fancyhf{}
\renewcommand{\headrulewidth}{0.4pt}\renewcommand{\footrulewidth}{0.4pt}
\fancyhead[L]{\small   \textit{Banaras Hindu University}}
\fancyhead[R]{\small   \textit{GLB-101: Elementary Physical \& Structural Geology}}
\fancyfoot[C]{\small Page  \thepage\ of \pageref{LastPage}}


\newlist{parts}{enumerate}{2}
\setlist[parts,1]{label=(\alph*),leftmargin=*,itemsep=5pt,topsep=3pt}

\newcommand{\pts}[1]{\hfill   \textit{[#1]}}


\begin{document}

\begin{center}
  {\large\bfseries BANARAS HINDU UNIVERSITY}\\[2pt]
  {\normalsize B.Sc. (Hons.) --- Semester I Examination, 2023-24}\\[6pt]
  \rule{0.8\linewidth}{0.5pt}\\[6pt]
  {\large\bfseries Paper: GLB-101 --- Elementary Physical and Structural Geology}\\[4pt]
  {\normalsize Time: 3 Hours \hfill Maximum Marks: 70}\\[2pt]
  {\small Ref. No.: 067572}
\end{center}
\rule{\linewidth}{0.4pt}
\smallskip



\noindent   \textit{Instructions: Answer   \textbf{five} questions in all, selecting at least   \textbf{two} each from Sections A and B, and Section-C which is compulsory.}
\bigskip

\bigskip


\noindent {\large\bfseries SECTION --- A}
\rule{\linewidth}{0.4pt}\smallskip




\noindent   \textbf{Question 1.} With a neat diagram of a fold, define different elements of it. Describe different types of fold based on the dip of the axial plane, plunge of the hinge line, and fold profile with neat diagrams. \pts{14}

\medskip


\noindent   \textbf{Question 2.} Describe the difference between the conformable and unconformable contacts in a sequence of strata. Describe the types of unconformities with neat diagrams. Discuss the importance of unconformities in geology. \pts{14}

\medskip


\noindent   \textbf{Question 3.} Discuss in detail the origin and evolution of the hydrosphere and biosphere. \pts{14}

\bigskip


\noindent {\large\bfseries SECTION --- B}
\rule{\linewidth}{0.4pt}\smallskip




\noindent   \textbf{Question 4.} Write short notes on the following: \pts{$2  \times 7 = 14$}

\begin{parts}
  \item Overlap, outlier and inlier
  \item Translational and rotational movement along faults with neat diagrams
\end{parts}

\medskip


\noindent   \textbf{Question 5.} Discuss with suitable examples: \pts{$2  \times 7 = 14$}

\begin{parts}
  \item Paleomagnetism and geomagnetic time scale
  \item Relative methods of age determination
\end{parts}

\medskip


\noindent   \textbf{Question 6.} Discuss the following: \pts{$2  \times 7 = 14$}

\begin{parts}
  \item Dip-slip fault and normal fault with diagrams
  \item Different types of chemical weathering processes
\end{parts}

\medskip


\noindent   \textbf{Question 7.} Briefly discuss the following: \pts{$4  \times 3\frac{1}{2} = 14$}

\begin{parts}
  \item Objectives of topographic maps
  \item Lopolith and phacolith
  \item Body fossil and trace fossil
  \item Crater and caldera
\end{parts}

\bigskip


\noindent {\large\bfseries SECTION --- C}
\rule{\linewidth}{0.4pt}\smallskip




\noindent   \textbf{Question 8.} Describe or fill in the blanks for each of the following in 2 to 3 sentences. Give suitable diagrams wherever necessary: \pts{14}

\begin{parts}
  \item Contour map of a plateau \pts{1}
  \item Types of scale in Structural Geology \pts{1}
  \item Importance of dip and strike \pts{1}
  \item Strike joint \pts{1}
  \item Upright and vertical folds \pts{1}
  \item Rake of slickenside present on a fault plane \pts{1}
  \item Azimuth circle on a clinometer compass \pts{1}
  \item The point of origin of an earthquake is known as \underline{\hspace{3cm}}. \pts{1}
  \item Among the commonly found radioisotopes, the longest half-life belongs to the \underline{\hspace{3cm}} decay series and the shortest half-life belongs to the \underline{\hspace{3cm}} decay series. \pts{2}
  \item Dinosaurs first appeared in the \underline{\hspace{3cm}} period and became extinct in the \underline{\hspace{3cm}} period. \pts{2}
  \item The protolith for mica-schist is \underline{\hspace{3cm}} and that for marble is \underline{\hspace{3cm}}. \pts{2}
\end{parts}

\vfill\rule{\linewidth}{0.4pt}

\begin{center}\small
\href{https://bhu-exam.vercel.app/}{Click here to visit our website}\\[4pt]
\href{https://me-aryan.vercel.app/}{Click here to visit our contributor}\\[4pt]
\href{https://quashan-me.vercel.app/}{Click here to visit our contributor}
\end{center}
\end{document}
